\section{AIME II 2002}
        \begin{enumerate}
        	\item (\textit{AIME II 2002}) Given that $x$ and $y$ are both integers between $100$ and $999$, inclusive; $y$ is the number formed by reversing the digits of $x$; and $z=|x-y|$. How many distinct values of $z$ are possible?


 
 \textbf{Answer: }009



\textbf{Solution}

We express the numbers as $x=100a+10b+c$ and $y=100c+10b+a$.  From this, we have 
 \begin{eqnarray*}z&=&|100a+10b+c-100c-10b-a|\\&=&|99a-99c|\\&=&99|a-c|\\ \end{eqnarray*}Because $a$ and $c$ are digits, and $a$ and $c$ are both between 1 and 9 (from condition 1), there are $\boxed{009}$ possible values (since all digits except $9$ can be expressed this way).

 


	\item (\textit{AIME II 2002}) Three \href{https://artofproblemsolving.com/wiki/index.php?title=Vertex}{vertices} of a \href{https://artofproblemsolving.com/wiki/index.php?title=Cube}{cube} are $P=(7,12,10)$, $Q=(8,8,1)$, and $R=(11,3,9)$. What is the \href{https://artofproblemsolving.com/wiki/index.php?title=Surface\_area}{surface area} of the cube?


 
 \textbf{Answer: }294



\textbf{Solution}

$PQ=\sqrt{(8-7)^2+(8-12)^2+(1-10)^2}=\sqrt{98}$
$PR=\sqrt{(11-7)^2+(3-12)^2+(9-10)^2}=\sqrt{98}$
$QR=\sqrt{(11-8)^2+(3-8)^2+(9-1)^2}=\sqrt{98}$

 So, $PQR$ is an equilateral triangle which means that $PQR$ are the diagonals of 3 sides that share a vertex. 

 Let the side of the cube be $a$.Since $PQ,PR,QR$ are diagonals $a\sqrt{2}=\sqrt{98}$

 So, $a=7$, and hence the surface area is $6a^2=\framebox{294}$.

 


	\item (\textit{AIME II 2002}) It is given that $\log_{6}a + \log_{6}b + \log_{6}c = 6$, where $a$, $b$, and $c$ are \href{https://artofproblemsolving.com/wiki/index.php?title=Positive}{positive} \href{https://artofproblemsolving.com/wiki/index.php?title=Integer}{integers} that form an increasing \href{https://artofproblemsolving.com/wiki/index.php?title=Geometric\_sequence}{geometric sequence} and $b - a$ is the \href{https://artofproblemsolving.com/wiki/index.php?title=Perfect\_square}{square} of an integer. Find $a + b + c$.


 
 \textbf{Answer: }111



\textbf{Solution 1}

$abc=6^6$. Since they form an increasing geometric sequence, $b$ is the \href{https://artofproblemsolving.com/wiki/index.php?title=Geometric\_mean}{geometric mean} of the \href{https://artofproblemsolving.com/wiki/index.php?title=Product}{product} $abc$. $b=\sqrt[3]{abc}=6^2=36$.

 Since $b-a$ is the square of an integer, we can find a few values of $a$ that work: $11, 20, 27, 32,$ and $35$. Out of these, the only value of $a$ that works is $a=27$, from which we can deduce that $c=\dfrac{36}{27}\cdot 36=\dfrac{4}{3}\cdot 36=48$.

 Thus, $a+b+c=27+36+48=\boxed{111}$




\textbf{Solution 2(similar to Solution 1)}

Let $r$ be the common ratio of the geometric sequence. Since it is increasing, that means that $b = ar$, and $c = ar^2$. Simplifying the logarithm, we get $\log_6(a^3 \cdot r^3) = 6$. Therefore, $a^3 \cdot r^3 = 6^6$. Taking the cube root of both sides, we see that $ar = 6^2 = 36$. Now since $ar = b$, that means $b = 36$. Using the trial and error shown in solution 1, we get $a = 27$, and $r = \frac{4}{3}$. Now, $27 \cdot r^2= c = 48$. Therefore, the answer is $27+36+48 = \boxed{111}$



	\item (\textit{AIME II 2002}) Patio blocks that are hexagons $1$ unit on a side are used to outline a garden by placing the blocks edge to edge with $n$ on each side. The diagram indicates the path of blocks around the garden when $n=5$.


 
                    \begin{figure}[H]
                        \centering
                        \includegraphics[width=0.5\textwidth]{Problem/AIME_II_2002_4_Problem_Figure_1.png}
                    \end{figure}
                    


 If $n=202$, then the area of the garden enclosed by the path, not including the path itself, is $m\left(\sqrt3/2\right)$ square units, where $m$ is a positive integer. Find the remainder when $m$ is divided by $1000$.


 
 \textbf{Answer: }803



\textbf{Solution 1}

When $n>1$, the path of blocks has $6(n-1)$ blocks total in it. When $n=1$, there is just one lonely block. Thus, the area of the garden enclosed by the path when $n=202$ is

 
 \[(1+6+12+18+\cdots +1200)A=(1+6(1+2+3...+200))A\],

 where $A$ is the area of one block. Then, because $n(n+1)/2$ is equal to the sum of the first $n$ integers:

 
 \[(1+6(1+2+3...+200))=(1+6((200)(201)/2))A=120601A\].

 Since $A=\dfrac{3\sqrt{3}}{2}$, the area of the garden is

 
 \[120601\cdot \dfrac{3\sqrt{3}}{2}=\dfrac{361803\sqrt{3}}{2}\].

 $m=361803$, $\dfrac{m}{1000}=361$ Remainder $\boxed{803}$.

 



\textbf{Solution 2}

Note that this is just the definition for a centered hexagonal number, and  the formula for $(n-1)^{th}$ term is $3n(n+1)+1$. Applying this for $200$ as we want the inner area gives $120601$. Then continue as above.

 



\textbf{Solution 3 (Some More Calculation, by DARSHSHETTY4)}

You can also find the area of the inside of the garden by constructing a larger hexagon using segments connected to the centers of each of the "vertex" mini-hexagons. 

 So, the side length of the larger constructed hexagon is (202-2) * √3 + √3/2 + √3/2 or 201 * √3. 

 So, the area of the larger constructed hexagon is (201*√3)\textasciicircum{}2 * 3 * √3/2. 

 Now, we need to find the area of the smaller hexagons that are between the garden boundary and the constructed hexagon boundary. The area of the side length's hexagons per side length are 202-2 * 3√3 / 2 / 2, or 200 * 3/2 * √3/2, and extrapolating it to the whole hexagon, it is 6 * 200 * 3/2 * √3/2. 

 Now, we need to account for the areas of the vertex hexagons. The segments create 90 degree angles on 2 sides of the hexagon, creating 2 30 60 90 right triangles (leg 1/2, √3/2 and 1 equilateral triangle (side length 1). 

 
So, the sum of the vertex hexagon areas are 6 * ((1\textasciicircum{}2 * √3/4) +√3/8 + √3/8). Hence, having the 2nd and 3rd expressions subtracted from the first, we get 361803 * √3/2, so the remainder is 803.

 


	\item (\textit{AIME II 2002}) Find the sum of all positive integers $a=2^n3^m$ where $n$ and $m$ are non-negative integers, for which $a^6$ is not a divisor of $6^a$.


 
 \textbf{Answer: }042



\textbf{Solution 1}

Substitute $a=2^n3^m$ into $a^6$ and $6^a$, and find all pairs of non-negative integers (n,m) for which $(2^n3^m)^{6}$ is not a divisor of $6^{2^n3^m}$

 Simplifying both expressions:

 $2^{6n} \cdot 3^{6m}$ is not a divisor of $2^{2^n3^m} \cdot 3^{2^n3^m}$

 Comparing both exponents (noting that there must be either extra powers of 2 or extra powers of 3 in the left expression):

 $6n > 2^n3^m$  OR  $6m > 2^n3^m$

 
Using the first inequality $6n > 2^n3^m$ and going case by case starting with n $\in$ {0, 1, 2, 3...}:

 n=0:$0>1 \cdot 3^m$ which has no solution for non-negative integers m 

 n=1:$6 > 2 \cdot 3^m$ which is true for m=0 but fails for higher integers $\Rightarrow (1,0)$

 n=2:$12 > 4 \cdot 3^m$ which is true for m=0 but fails for higher integers $\Rightarrow (2,0)$

 n=3:$18 > 8 \cdot 3^m$ which is true for m=0 but fails for higher integers $\Rightarrow (3,0)$

 n=4:$24 > 16 \cdot 3^m$ which is true for m=0 but fails for higher integers $\Rightarrow (4,0)$

 n=5:$30 > 32 \cdot 3^m$ which has no solution for non-negative integers m 

 There are no more solutions for higher $n$, as polynomials like $6n$ grow slower than exponentials like $2^n$.

 


 Using the second inequality $6m > 2^n3^m$ and going case by case starting with m $\in$ {0, 1, 2, 3...}:

 m=0:$0>2^n \cdot 1$ which has no solution for non-negative integers n 

 m=1:$6>2^n \cdot 3$ which is true for n=0 but fails for higher integers $\Rightarrow (0,1)$

 m=2:$12>2^n \cdot 9$ which is true for n=0 but fails for higher integers $\Rightarrow (0,2)$

 m=3:$18>2^n \cdot 27$ which has no solution for non-negative integers n 

 There are no more solutions for higher $m$, as polynomials like $6m$ grow slower than exponentials like $3^m$.

 
Thus there are six numbers corresponding to (1,0), (2,0), (3,0), (4,0), (0,1), and (0,2).  Plugging them back into the original expression, these numbers are 2, 4, 8, 16, 3, and 9, respectively.  Their sum is $\framebox{042}$.

 



\textbf{Solution 2 (faster and more concise)}

Notice that the condition is equivalent to saying

 
 \[v_2(a^6) \geq v_2(6^a) \implies 6n \geq a\]
 \[v_3(a^6) \geq v_3(6^a) \implies 6m \geq a.\]

 Notice that we cannot have both expressions to be equality state, as that would result in $a^6 = 6^a.$ Testing, we see the possible pairs $(n, m)$ are $(1, 0), (2, 0), (3, 0), (4, 0), (0, 1), (0, 2).$  Since the $a$ grows much faster than the left-hand side of the above inequalities, these are all possible solutions.  Adding, we get $\framebox{042}$. 

 


	\item (\textit{AIME II 2002}) Find the integer that is closest to $1000\sum_{n=3}^{10000}\frac1{n^2-4}$.


 
 \textbf{Answer: }521



\textbf{Solution 2}

Using the fact that $\frac{1}{n(n+k)} = \frac{1}{k} ( \frac{1}{n}-\frac{1}{n+k} )$ or by partial fraction decomposition, we both obtained $\frac{1}{x^2-4} = \frac{1}{4}(\frac{1}{x-2}-\frac{1}{x+2})$. The denominators of the positive terms are $1,2,..,9998$, while the negative ones are $5,6,...,10002$. Hence we are left with $1000 \cdot \frac{1}{4} (1 + \frac{1}{2} + \frac{1}{3} + \frac{1}{4} - \frac{1}{9999} - \frac{1}{10000} - \frac{1}{10001} - \frac{1}{10002})$. We can simply ignore the last $4$ terms, and we get it is approximately $1000*\frac{25}{48}$. Computing $\frac{25}{48}$ which is about $0.5208..$ and moving the decimal point three times, we get that the answer is $521$



	\item (\textit{AIME II 2002}) It is known that, for all positive integers $k$,


 $1^2+2^2+3^2+\ldots+k^{2}=\frac{k(k+1)(2k+1)}6$.

 Find the smallest positive integer $k$ such that $1^2+2^2+3^2+\ldots+k^2$ is a multiple of $200$.


 
 \textbf{Answer: }112



\textbf{Solution}

Using the given formula, we require $\frac{k(k+1)(2k+1)}{6}$ to be a multiple of $200$, i.e. $k(k+1)(2k+1)$ to be a multiple of $200 \cdot 6 = 1200 = 2^4 \cdot 3 \cdot 5^2$. This is equivalent to $k(k+1)(2k+1)$ being divisible by each of these factors ($2^4 = 16$, $3$, and $5^2 = 25$) separately, since they are coprime.

 Thus, for divisibility by $16$, we observe that $(2k+1)$ is necessarily odd, while exactly one of $k$ and $(k+1)$ is even. It follows that the (at least) $4$ factors of $2$ must be either all contained in $k$, or all contained in $(k+1)$, yielding $k \equiv 0 \pmod{16}$ or $k+1 \equiv 0 \pmod{16}$ respectively. This reduces to $k \in \left\{0,15\right\} \pmod{16}$.

 Next, if $k \equiv 0 \pmod{3}$, then obviously $k(k+1)(2k+1)$ will be divisible by 3; otherwise, we will have either $k \equiv 1 \pmod{3}$, giving $2k+1 \equiv 2 \cdot 1 + 1 \equiv 3 \equiv 0 \pmod{3}$, or $k \equiv 2 \pmod{3}$, giving $k+1 \equiv 2+1 \equiv 3 \equiv 0 \pmod{3}$. Hence $k(k+1)(2k+1)$ is in fact always divisible by $3$, regardless of the value of $k$.

 Similarly, considering the cases $k \equiv 0,1,2,3,4 \pmod{5}$ in turn, the residues modulo $5$ of $k$, $(k+1)$, and $(2k+1)$ respectively are $0,1,1$; $1,2,3$; $2,3,0$; $3,4,2$; and $4,0,4$. As $0$ never appears more than once in any of these cases, we deduce that at most one of $k$, $(k+1)$, and $(2k+1)$ is divisible by 5. Analogously to above, it follows that the (at least) $2$ factors of $5$ must be either all contained in $k$, all contained in $(k+1)$, or all contained in $(2k+1)$. These respectively give $k,k+1,2k+1 \equiv 0 \pmod{25}$, which reduces to $k \in \left\{0,12,24\right\} \pmod{25}$. 

 Accordingly, as there are $2$ possible residues for $k$ modulo $16$ and $3$ possible residues modulo $25$, we obtain a total of $2 \cdot 3 = 6$ pairs of simultaneous congruences. By the \href{https://artofproblemsolving.com/wiki/index.php?title=Chinese\_Remainder\_Theorem}{Chinese remainder theorem}, as $16$ and $25$ are coprime, each pair has a unique solution modulo $16 \cdot 25 = 400$, which we can easily calculate by simply writing out the solutions of the first congruence, then checking whether each of them also satisfies the second congruence.

 For instance, to solve $k \equiv 0 \pmod{16}$, $k \equiv 12 \pmod{25}$, we can write out the positive multiples of $16$ (to satisfy the first congruence, together with the fact that $k > 0$), then subtract $12$ from each, until we find a multiple of $16$ for which this subtraction gives a multiple of $25$. As it turns out, $16 \cdot 7 = 112$ and $112-12 = 100 = 25 \cdot 4$, so the solution of this pair is precisely $k \equiv 112 \pmod{400}$. By applying this algorithm to each of the remaining pairs of congruences, we eventually obtain the complete solution $k \in \left\{0,112,175,224,287,399\right\} \pmod{400}$, so the smallest possible positive value of $k$ is $\boxed{112}$.

 


	\item (\textit{AIME II 2002}) Find the least positive integer $k$ for which the equation $\left\lfloor\frac{2002}{n}\right\rfloor=k$ has no integer solutions for $n$. (The notation $\lfloor x\rfloor$ means the greatest integer less than or equal to $x$.)


 
 \textbf{Answer: }049



\textbf{Solution 1}

Note that if $\frac{2002}n - \frac{2002}{n+1}\leq 1$, then either $\left\lfloor\frac{2002}{n}\right\rfloor=\left\lfloor\frac{2002}{n+1}\right\rfloor$,or $\left\lfloor\frac{2002}{n}\right\rfloor=\left\lfloor\frac{2002}{n+1}\right\rfloor+1$. Either way, we won't skip any natural numbers.

 The greatest $n$ such that $\frac{2002}n - \frac{2002}{n+1} > 1$ is $n=44$. (The inequality simplifies to $n(n+1)<2002$, which is easy to solve by trial, as the solution is obviously $\simeq \sqrt{2002}$.)

 
We can now compute:
 \[\left\lfloor\frac{2002}{45}\right\rfloor=44\]
 \[\left\lfloor\frac{2002}{44}\right\rfloor=45\]
 \[\left\lfloor\frac{2002}{43}\right\rfloor=46\]
 \[\left\lfloor\frac{2002}{42}\right\rfloor=47\]
 \[\left\lfloor\frac{2002}{41}\right\rfloor=48\]
 \[\left\lfloor\frac{2002}{40}\right\rfloor=50\]

 From the observation above (and the fact that $\left\lfloor\frac{2002}{2002}\right\rfloor=1$) we know that all integers between $1$ and $44$ will be achieved for some values of $n$. Similarly, for $n<40$ we obviously have $\left\lfloor\frac{2002}{n}\right\rfloor > 50$.

 Hence the least positive integer $k$ for which the equation $\left\lfloor\frac{2002}{n}\right\rfloor=k$ has no integer solutions for $n$ is $\boxed{049}$.

 \textbf{Note} After getting that $\left\lfloor\frac{2002}{45}\right\rfloor=44$, for ease of computation above, we can use the fact that $(40+k)(49-k)$ varies solely based on $k^2$ and checking these gives us that the pattern fails at $k=0$ giving us $\boxed{049}$ as the answer.

 



\textbf{Solution 2}

Rewriting the given information and simplifying it a bit, we have
 \begin{align*}  k \le \frac{2002}{n} < k+1 &\implies \frac{1}{k} \ge \frac{n}{2002} > \frac{1}{k+1}. \\ &\implies \frac{2002}{k} \ge n > \frac{2002}{k+1}.   \end{align*}


 Now note that in order for there to be no integer solutions to $n,$ we must have $\left\lfloor \frac{2002}{k} \right\rfloor = \left\lfloor \frac{2002}{k+1} \right\rfloor.$ We seek the smallest such $k.$ A bit of experimentation yields that $k=49$ is the smallest solution, as for $k=49,$ it is true that $\left\lfloor \frac{2002}{49} \right\rfloor = \left\lfloor \frac{2002}{50} \right\rfloor = 40.$ Furthermore, $k=49$ is the smallest such case. (If unsure, we could check if the result holds for $k=48,$ and as it turns out, it doesn't.) Therefore, the answer is $\boxed{049}.$




\textbf{Solution 3}

In this solution we use inductive reasoning and a lot of trial and error. Depending on how accurately you can estimate, the solution will come quicker or slower.

 Using values of $k$ as $1, 2, 3, 4,$ and $5,$ we can find the corresponding values of $n$ relatively easily. For $k = 1$, $n$ is in the range $[2002-1002]$; for $k = 2$, $n$ is the the range $[1001-668]$, etc: $3, [667,501]; 4, [500-401]; 5, [400-334]$. For any positive integer $k, n$ is in a range of $\left\lfloor \frac{2002}{k} \right\rfloor -\left\lceil \frac{2002}{k+1} \right\rceil$.

 Now we try testing $k = 1002$ to get a better understanding of what our solution will look like. Obviously, there will be no solution for $n$, but we are more interested in how the range will compute to. Using the formula we got above, the range will be $1-2$. Testing any integer $k$ from $1002-2000$ will result in the same range. Also, notice that each and every one of them have no solution for $n$. Testing $1001$ gives a range of $2-2$, and $2002$ gives $1-1$. They each have a solution for $n$, and their range is only one value. Therefore, we can assume with relative safety that the integer $k$ we want is the lowest integer that follows this equation

 
 \[\left\lfloor\frac{2002}{k}\right\rfloor + 1 = \left\lceil \frac{2002}{k+1}\right\rceil\]

 Now we can easily guess and check starting from $k = 1$. After a few tests it's not difficult to estimate a few jumps, and it took me only a few minutes to realize the answer was somewhere in the forties (You could also use the fact that $45^2=2025$). Then it's just a matter of checking them until we get $\boxed{049}$.Alternatively, you could use the equation above and proceed with one of the other two solutions listed.

 
-jackshi2006

 Edited and $\LaTeX$ed by PhunsukhWangdu

 



\textbf{Solution 4}

Here is an intuitive way to approximate the answer is around $45$: For the function $f(n)=\frac{2002}{n}$ its derivative is $-\frac{2002}{n^2}$,which should be close to $-1$ because we need to find the smallest skipped integer. The rest of the steps are the same as Solution 1.

 -maxamc

 


	\item (\textit{AIME II 2002}) Let $\mathcal{S}$ be the set $\lbrace1,2,3,\ldots,10\rbrace$. Let $n$ be the number of sets of two non-empty disjoint subsets of $\mathcal{S}$. (Disjoint sets are defined as sets that have no common elements.) Find the remainder obtained when $n$ is divided by $1000$.


 
 \textbf{Answer: }501



\textbf{Solution 1 (combinatorial)}

Let the two disjoint subsets be $A$ and $B$, and define $C = \mathcal{S}\setminus\left(A \cup B\right)$ (so $C$ is disjoint from both $A$ and $B$, and every element of $\mathcal{S}$, if not contained in $A$ or $B$, is necessarily contained in $C$). This means every element $s \in \mathcal{S}$ satisfies exactly one of $s \in A$, $s \in B$, or $s \in C$, i.e. for each of the $10$ elements, we have to independently choose $1$ of these $3$ subsets into which to place it. Hence there are a total of $3^{10}$ ways to partition $\mathcal{S}$ into disjoint subsets $A$, $B$, and $C$. 

 However, we must exclude those partitions in which $A$ or $B$ (or both) is empty. If $A = \emptyset$, and thus 

 


	\item (\textit{AIME II 2002}) While finding the sine of a certain angle, an absent-minded professor failed to notice that his calculator was not in the correct angular mode. He was lucky to get the right answer. The two least positive real values of $x$ for which the sine of $x$ degrees is the same as the sine of $x$ radians are $\frac{m\pi}{n-\pi}$ and $\frac{p\pi}{q+\pi}$, where $m$, $n$, $p$, and $q$ are positive integers. Find $m+n+p+q$.


 
 \textbf{Answer: }900



\textbf{Solution 1}

Note that $x$ degrees is equal to $\frac{\pi x}{180}$ radians. Also, for $\alpha \in \left[0 , \frac{\pi}{2} \right]$, the two least positive angles $\theta > \alpha$ such that $\sin{\theta} = \sin{\alpha}$ are $\theta = \pi-\alpha$, and $\theta = 2\pi + \alpha$. 

 Clearly $x > \frac{\pi x}{180}$ for positive real values of $x$. 

 $\theta = \pi-\alpha$ yields: $x = \pi - \frac{\pi x}{180} \Rightarrow x = \frac{180\pi}{180+\pi} \Rightarrow (p,q) = (180,180)$. 

 $\theta = 2\pi + \alpha$ yields: $x = 2\pi + \frac{\pi x}{180} \Rightarrow x = \frac{360\pi}{180-\pi} \Rightarrow (m,n) = (360,180)$. 

 So, $m+n+p+q = \boxed{900}$.

 



\textbf{Solution 2}

The first case is when the two angles, $x$ and $\frac{\pi x}{180}$, are coterminal. The second case is when they are reflections of the $y$ axis.

 1. $x+2\pi a = \frac{\pi x}{180}$ for any integer $a$So $x=\frac{360\pi a }{\pi -180}$

 2. $(2b+1)\pi -x = \frac{\pi x}{180}$ for any integer $b$So $x = \frac{180(2b+1)\pi}{\pi + 180}$

 Choosing carefully $a,b$ such that it's the minimum gives the answer same as above.

 


	\item (\textit{AIME II 2002}) Two distinct, real, infinite geometric series each have a sum of $1$ and have the same second term. The third term of one of the series is $1/8$, and the second term of both series can be written in the form $\frac{\sqrt{m}-n}p$, where $m$, $n$, and $p$ are positive integers and $m$ is not divisible by the square of any prime. Find $100m+10n+p$.


 
 \textbf{Answer: }518



\textbf{Solution 1}

Let the second term of each series be $x$. Then, the common ratio is $\frac{1}{8x}$, and the first term is $8x^2$. 

 So, the sum is $\frac{8x^2}{1-\frac{1}{8x}}=1$. Thus, $64x^3-8x+1 = (4x-1)(16x^2+4x-1) = 0 \Rightarrow x = \frac{1}{4}, \frac{-1 \pm \sqrt{5}}{8}$. 

 The only solution in the appropriate form is $x = \frac{\sqrt{5}-1}{8}$. Therefore, $100m+10n+p = \boxed{518}$.

 


 



\textbf{Solution 2}

Let the two sequences be $a, ar, ar^2 ... \text{ }an^2$ and $x, xy, xy^2 ... \text{ }xy^n$. We know for a fact that $ar = xy$. We also know that the sum of the first sequence = $\frac{a}{1-r} = 1$, and the sum of the second sequence = $\frac{x}{1-y} = 1$. Therefore we have 
 \[a+r = 1\]
 \[x+y = 1\]
 \[ar=xy\] We can then replace $r = \frac{xy}{a}$ and $y = \frac{ar}{x}$. We plug them into the two equations $a+r = 1$ and $x+y = 1$. We then get 
 \[x^2 + ar = x\]
 \[a^2 + xy = a\]We subtract these equations, getting 
 \[x^2 - a^2 + ar - xy = x-a\]Remember 
 \[ar=xy\], so 
 \[(x-a)(x+a-1) = 0\]Then considering cases, we have either $x=a$ or $y=a$. This suggests that the second sequence is in the form $r, ra, ra^2...$, while the first sequence is in the form $a, ar, ar^2...$ Now we have that $ar^2 = \frac18$ and we also have that $a+r = 1$. We can solve for $r$ and the only appropriate value for $r$ is $\frac{1+\sqrt{5}}{4}$. All we want is the second term, which is $ar = \frac{ar^2}{r} = \frac{\frac18}{\frac{1+\sqrt{5}}{4}} = \frac{\sqrt{5} - 1}{8}$solution by jj\_ca888

 



\textbf{Solution 3}

Let's ignore the "two distinct, real, infinite geometric series" part for now and focus on what it means to be a geometric series.

 Let the first term of the series with the third term equal to $\frac18$ be $a,$ and the common ratio be $r.$ Then, we get that $\frac{a}{1-r} = 1 \implies a = 1-r,$ and $ar^2 = \frac18 \implies (1-r)(r^2) = \frac18.$

 We see that this cubic is equivalent to $r^3 - r^2 + \frac18 = 0.$ Through experimenting, we find that one of the solutions is $r = \frac12.$ Using synthetic division leads to the quadratic $4x^2 - 2x - 1 = 0.$ This has roots $\dfrac{2 \pm \sqrt{4 - 4(4)(-1)}}{8},$ or, when reduced, $\dfrac{1 \pm \sqrt{5}}{4}.$

 It becomes clear that the two geometric series have common ratio $\frac{1 + \sqrt{5}}{4}$ and $\frac{1 - \sqrt{5}}{4}.$ Let $\frac{1 + \sqrt{5}}{4}$ be the ratio that we are inspecting. We see that in this case, $a = \dfrac{3 - \sqrt{5}}{4}.$

 Since the second term in the series is $ar,$ we compute this and have that 
 \[ar = \left(\dfrac{3 - \sqrt{5}}{4} \right)\left(\dfrac{1+\sqrt{5}}{4}\right) = \dfrac{\sqrt{5} - 1}{8},\]for our answer of $100 \cdot 5 + 1 \cdot 10 + 8 = \boxed{518}.$

 Solution by Ilikeapos

 \textbf{Sidenote} One of the geometric series has first term $\frac{3 - \sqrt{5}}{4}$ and common ratio $\frac{1 + \sqrt{5}}{4}$, and its third term is $\frac{1}{8}$. The other series has first term $\frac{1 + \sqrt{5}}{4}$ and common ratio $\frac{3 - \sqrt{5}}{4}$.

 


	\item (\textit{AIME II 2002}) A basketball player has a constant probability of $.4$ of making any given shot, independent of previous shots. Let $a_n$ be the ratio of shots made to shots attempted after $n$ shots. The probability that $a_{10}=.4$ and $a_n\le.4$ for all $n$ such that $1\le n\le9$ is given to be $p^aq^br/\left(s^c\right)$ where $p$, $q$, $r$, and $s$ are primes, and $a$, $b$, and $c$ are positive integers. Find $\left(p+q+r+s\right)\left(a+b+c\right)$.


 
 \textbf{Answer: }660



\textbf{Solution 1}

We graph the $10$ shots on a grid. Suppose that a made shot is represented by a step of $(0,1)$, and a missed shot is represented by $(1,0)$. Then the basketball player's shots can be represented by the number of paths from $(0,0)$ to $(6,4)$ that always stay below the line $y=\frac{2x}{3}$. We can find the number of such paths using a Pascal's Triangle type method below, computing the number of paths to each point that only move right and up.
                    \begin{figure}[H]
                        \centering
                        \includegraphics[width=0.5\textwidth]{Solution/AIME_II_2002_12_Solution_Figure_1.png}
                    \end{figure}
                    Therefore, there are $23$ ways to shoot $4$ makes and $6$ misses under the given conditions. The probability of each possible sequence occurring is $(.4)^4(.6)^6$. Hence the desired probability is
 \[\frac{23\cdot 2^4\cdot 3^6}{5^{10}},\]and the answer is $(23+2+3+5)(4+6+10)=\boxed{660}$.

 



\textbf{Solution 2}

The first restriction is that $a_{10} = .4$, meaning that the player gets exactly 4 out of 10 baskets.  The second restriction is $a_n\le.4$.  This means that the player may never have a shooting average over 40\%.  Thus, the first and second shots must fail, since $\frac{1}{1}$ and $\frac{1}{2}$ are both over $.4$, but the player may make the third basket, since $\frac{1}{3} \le .4$  In other words, the earliest the first basket may be made is attempt 3.  Using similar reasoning, the earliest the second basket may be made is attempt 5, the earliest the third basket may be made is attempt 8, and the earliest the fourth basket may be made is attempt 10.

 Using X to represent a basket and O to represent a failure, this 'earliest' solution may be represented as:

 OOXOXOOXOX

 To simplify counting, note that the first, second, and tenth shots are predetermined.  The first two shots must fail, and the last shot must succeed.  Thus, only slots 3-9 need to be counted, and can be abbreviated as follows:

 XOXOOXO

 The problem may be separated into five cases, since the first shot may be made on attempt 3, 4, 5, 6, or 7.  The easiest way to count the problem is to remember that each X may slide to the right, but NOT to the left.

 First shot made on attempt 3:

 XOXOOXO

 XOXOOOX

 XOOXOXO

 XOOXOOX

 XOOOXXO

 XOOOXOX

 XOOOOXX

 \textbf{Total - 7}

 First shot made on attempt 4:

 Note that all that needs to be done is change each line in the prior case from starting with "XO....." to "OX.....".

 \textbf{Total - 7}

 First shot made on attempt 5:

 OOXXOXO

 OOXXOOX

 OOXOXXO

 OOXOXOX

 OOXOOXX

 \textbf{Total - 5}

 First shot made on attempt 6:

 OOOXXXO

 OOOXXOX

 OOOXOXX

 \textbf{Total - 3}

 First shot made on attempt 7:

 OOOOXXX

 \textbf{Total - 1}

 The total number of ways the player may satisfy the requirements is $7+7+5+3+1=23$.

 The chance of hitting any individual combination (say, for example, OOOOOOXXXX) is $\left(\frac{3}{5}\right)^6\left(\frac{2}{5}\right)^4$

 Thus, the chance of hitting any of these 23 combinations is $23\left(\frac{3}{5}\right)^6\left(\frac{2}{5}\right)^4 = \frac{23\cdot3^6\cdot2^4}{5^{10}}$

 Thus, the final answer is $(23+3+2+5)(6+4+10)=\boxed{660}$




\textbf{Solution 3}

Note $a_{10}=.4$. Therefore the player made 4 shots out of 10. He must make the 10th shot, because if he doesn't, then $a_9=\frac{4}{9}>.4$. Since $a_n\leq .4$ for all $n$ less than 11, we know that $a_1=a_2=0$. Now we must look at the 3rd through 9th shot.

 Now let's take a look at those un-determined shots. Let's put them into groups: the 3rd, 4th, and 5th shots in group A, and the 6th, 7th, 8th, and 9th shots in group B. The total number of shots made in groups A and B must be 3, since the player makes the 10th shot. We cannot have all three shots made in group A, since $a_5\leq .4$. Therefore we can have two shots made, one shot made, or no shots made in group A.

 \textbf{Case 1: Group A contains no made shots.}

 The number of ways this can happen in group A is 1. Now we must arrange the shots in group B accordingly. There are four ways to arrange them total, and all of them work. There are $\textbf{4}$ possibilities here.

 \textbf{Case 2: Group A contains one made shot.}

 The number of ways this could happen in group A is 3.  Now we must arrange the shots in group B accordingly. There are six ways to arrange them total, but the arrangement "hit hit miss miss" fails, because that would mean $a_7=\frac{3}{7}>.4$. All the rest work. Therefore there are $3\cdot5=\textbf{15}$ possibilities here.

 \textbf{Case 3: Group A contains two made shots.}

 The number of ways this could happen in group A is 2 (hit hit miss doesn't work but the rest do). Now we must arrange the shots in group B accordingly. Note hit miss miss miss and miss hit miss miss fail. Therefore there are only 2 ways to do this, and there are $2\cdot 2=\textbf{4}$ total possibilities for this case.

 
Taking all these cases into account, we find that there are $4+15+4=23$ ways to have $a_{10} = .4$ and $a_n\leq .4$. Each of these has a probability of $.4^4\cdot.6^6=\frac{2^4\cdot 3^6}{5^{10}}$. Therefore the probability that we have $a_{10} = .4$ and $a_n\leq .4$ is $\frac{23\cdot2^4\cdot3^6}{5^{10}}$. Now we are asked to find the product of the sum of the primes and the sum of the exponents, which is $(23+2+3+5)(4+6+10)=33\cdot20=\boxed{660}$.

 


 



\textbf{Solution 4}

To satisfy $a_{10} = 0.4$, the player must make exactly 4 shots out of 10. Let $k_n$ be the total number of shots made after $n$ attempts. The condition $a_n \le 0.4$ means $k_n \le 0.4n$. The maximum allowable values of $k_n$ are:

 
 \[\begin{array}{|c|c|c|c|c|c|c|c|c|c|c|} \hline \text{Shot } (n) & 1 & 2 & 3 & 4 & 5 & 6 & 7 & 8 & 9 & 10 \\ \hline \text{Max Makes } (k_n) & 0 & 0 & 1 & 1 & 2 & 2 & 2 & 3 & 3 & 4 \\ \hline \end{array}\]

 At $n=9$, the player can have at most 3 makes. Since 4 makes are required by $n=10$, the tenth shot must be a make. Split the sequence into Group A (shots 1–7) and Group B (shots 8–10). The remaining three makes must be distributed between these two groups.

 \textbf{Case 1: 1 shot made in Group A, 3 shots made in Group B}

 Group A: The make cannot be on shots 1 or 2, since the maximum allowed number of makes there is 0. Thus the make may occur on any of shots 3 through 7. There are 5 ways.

 Group B: All three shots must be makes. At $n=9$, $k_9 = 3$, and $3/9 \le 0.4$, so this case is valid. There is 1 way.

 Subtotal: $5 \times 1 = 5$ ways.

 \textbf{Case 2: 2 shots made in Group A, 2 shots made in Group B}

 Group A: Two makes may be placed as follows.If both occur among shots 5, 6, or 7, there are $\binom{3}{2} = 3$ ways.If one occurs among shots 3 or 4 and one among shots 5, 6, or 7, there are $2 \times 3 = 6$ ways.Thus Group A contributes $3 + 6 = 9$ ways.

 Group B: Since shot 10 must be a make, the remaining make must occur on shot 8 or 9. There are 2 ways.

 Subtotal: $9 \times 2 = 18$ ways.

 \textbf{Final Calculation}

 The total number of valid sequences is $5 + 18 = 23$. Each sequence has probability
 \[(0.4)^4(0.6)^6 = \left(\frac{2}{5}\right)^4\left(\frac{3}{5}\right)^6 = \frac{2^4 \cdot 3^6}{5^{10}}.\]Thus the desired probability is
 \[\frac{23 \cdot 2^4 \cdot 3^6}{5^{10}}.\]

 Writing this in the form $\frac{p^a q^b r}{s^c}$ gives $p=2$, $q=3$, $r=23$, $s=5$, with $a=4$, $b=6$, and $c=10$. Therefore,
 \[(p+q+r+s)(a+b+c) = 33 \cdot 20 = \boxed{660}.\]



	\item (\textit{AIME II 2002}) In triangle $ABC$, point $D$ is on $\overline{BC}$ with $CD=2$ and $DB=5$, point $E$ is on $\overline{AC}$ with $CE=1$ and $EA=3$, $AB=8$, and $\overline{AD}$ and $\overline{BE}$ intersect at $P$. Points $Q$ and $R$ lie on $\overline{AB}$ so that $\overline{PQ}$ is parallel to $\overline{CA}$ and $\overline{PR}$ is parallel to $\overline{CB}$. It is given that the ratio of the area of triangle $PQR$ to the area of triangle $ABC$ is $m/n$, where $m$ and $n$ are relatively prime positive integers. Find $m+n$.


 
 \textbf{Answer: }901



\textbf{Solution 1}

Let $X$ be the intersection of $\overline{CP}$ and $\overline{AB}$. 

 
                    \begin{figure}[H]
                        \centering
                        \includegraphics[width=0.5\textwidth]{Solution/AIME_II_2002_13_Solution_Figure_1.png}
                    \end{figure}
                    

 Since $\overline{PQ} \parallel \overline{CA}$ and $\overline{PR} \parallel \overline{CB}$, $\angle CAB = \angle PQR$ and $\angle CBA = \angle PRQ$. So $\Delta ABC \sim \Delta QRP$, and thus, $\frac{[\Delta PQR]}{[\Delta ABC]} = \left(\frac{PX}{CX}\right)^2$. 

 Using \href{https://artofproblemsolving.com/wiki/index.php?title=Mass\_points}{mass points}: 

 \href{https://artofproblemsolving.com/wiki/index.php?title=WLOG}{WLOG}, let $W_C=15$. 

 Then: 

 $W_A = \left(\frac{CE}{AE}\right)W_C = \frac{1}{3}\cdot15=5$.

 $W_B = \left(\frac{CD}{BD}\right)W_C = \frac{2}{5}\cdot15=6$. 

 $W_X=W_A+W_B=5+6=11$. 

 $W_P=W_C+W_X=15+11=26$. 

 Thus, $\frac{PX}{CX}=\frac{W_C}{W_P}=\frac{15}{26}$. Therefore, $\frac{[\Delta PQR]}{[\Delta ABC]} = \left( \frac{15}{26} \right)^2 = \frac{225}{676}$, and $m+n=\boxed{901}$. Note we can just use mass points to get  $\left( \frac{15}{26} \right)^2= \frac{225}{676}$ which is $\boxed{901}$.

 



\textbf{Solution 2}

First draw $\overline{CP}$ and extend it so that it meets with $\overline{AB}$ at point $X$.

 
                    \begin{figure}[H]
                        \centering
                        \includegraphics[width=0.5\textwidth]{Solution/AIME_II_2002_13_Solution_Figure_1.png}
                    \end{figure}
                    

 We have that $[ABC]=\frac{1}{2}\cdot AC \cdot BC\sin{C}=\frac{1}{2}\cdot 4\cdot {7}\sin{C}=14\sin{C}$

 By Ceva's, 
 \[3\cdot{\frac{2}{5}}\cdot{\frac{BX}{AX}}=1\implies BX=\frac{5\cdot AX}{6}\] That means that 
 \[\frac{11\cdot {AX}}{6}=8\implies AX=\frac{48}{11} \ \text{and} \ BX=\frac{40}{11}\]

 Now we apply mass points. Assume WLOG that $W_{A}=1$. That means that

 
 \[W_{C}=3, W_{B}=\frac{6}{5}, W_{X}=\frac{11}{5}, W_{D}=\frac{21}{5}, W_{E}=4, W_{P}=\frac{26}{5}\]

 Notice now that $\triangle{PBQ}$ is similar to $\triangle{EBA}$. Therefore,

 
 \[\frac{PQ}{EA}=\frac{PB}{EB}\implies \frac{PQ}{3}=\frac{10}{13}\implies PQ=\frac{30}{13}\]

 Also, $\triangle{PRA}$ is similar to $\triangle{DBA}$. Therefore,

 
 \[\frac{PA}{DA}=\frac{PR}{DB}\implies \frac{21}{26}=\frac{PR}{5}\implies PR=\frac{105}{26}\]

 Because $\triangle{PQR}$ is similar to $\triangle{CAB}$, $\angle{C}=\angle{P}$.

 As a result, $[PQR]=\frac{1}{2}\cdot PQ \cdot PR \sin{C}=\frac{1}{2}\cdot \frac{30}{13}\cdot \frac{105}{26}\sin{P}=\frac{1575}{338}\sin{C}$.

 Therefore, 
 \[\frac{[PQR]}{[ABC]}=\frac{\frac{1575}{338}\sin{C}}{14\sin{C}}=\frac{225}{676}\implies 225+676=\boxed{901}\]

            \begin{itemize}
        \item Not the author writing here, but a note is that Ceva's Theorem was actually not necessary to solve this problem. The information was just nice to know :)
            \end{itemize}
            



\textbf{Solution 3}


                    \begin{figure}[H]
                        \centering
                        \includegraphics[width=0.5\textwidth]{Solution/AIME_II_2002_13_Solution_Figure_3.png}
                    \end{figure}
                    Use the mass of point. Denoting the mass of $C=15,B=6,A=5,D=21,E=20$, we can see that the mass of $P$ is $26$, hence we know that $\frac{BP}{PE}=\frac{10}{3}$, now we can find that $\frac{PQ}{AE}=\frac{10}{13}$ which implies $PQ=\frac{30}{13}$, it is obvious that $\triangle{PQR}$ is similar to $\triangle{ACB}$ so we need to find the ration between PQ and AC, which is easy, it is $\frac{15}{26}$, so our final answer is $\left( \frac{15}{26} \right)^2= \frac{225}{676}$ which is $\boxed{901}$. ~bluesoul

 



\textbf{Solution 4(Ceva & Menelaus)}

Construct $\overline{CP}$ and extend it to line $\overline{AB}$ at point $F$.

 
                    \begin{figure}[H]
                        \centering
                        \includegraphics[width=0.5\textwidth]{Solution/AIME_II_2002_13_Solution_Figure_4.png}
                    \end{figure}
                    

 Using Ceva's Theorem on triangle $ABC$ and point $P$, we get 
 \[\frac{AF}{BF} \times \frac{BD}{CD} \times \frac{CE}{AE}=1\]Thus, $\frac{AF}{BF}=\frac{6}{5}$.Then, using this info, we apply Menelaus on triangle $ACF$ and line $BE$, obtaining 
 \[\frac{AE}{BE} \times \frac{CP}{FP} \times \frac{FB}{AB}=1\]Simplifying and substituting, we find that $\frac{CP}{FP}=\frac{11}{15}$.Alternatively, $\frac{FP}{CF}=\frac{15}{26}$, which is also the ratio between the heights of the desired triangles.Finishing, $\left( \frac{15}{26} \right)^2= \frac{225}{676}$ achieving the final answer of $\boxed{901}$.~faliure167

 


	\item (\textit{AIME II 2002}) The perimeter of triangle $APM$ is $152$, and the angle $PAM$ is a right angle. A circle of radius $19$ with center $O$ on $\overline{AP}$ is drawn so that it is tangent to $\overline{AM}$ and $\overline{PM}$. Given that $OP=m/n$ where $m$ and $n$ are relatively prime positive integers, find $m+n$.


 
 \textbf{Answer: }098



\textbf{Solution 1 heheh}

Let the circle intersect $\overline{PM}$ at $B$. Then note $\triangle OPB$ and $\triangle MPA$ are similar. Also note that $AM = BM$ by \href{https://artofproblemsolving.com/wiki/index.php?title=Power\_of\_a\_point}{power of a point}. Using the fact that the ratio of corresponding sides in similar triangles is equal to the ratio of their perimeters, we have
 \[\frac{19}{AM} = \frac{152-2AM-19+19}{152} = \frac{152-2AM}{152}\]Solving, $AM = 38$. So the ratio of the side lengths of the triangles is 2. Therefore,
 \[\frac{PB+38}{OP}= 2 \text{ and } \frac{OP+19}{PB} = 2\]so $2OP = PB+38$ and $2PB = OP+19.$ Substituting for $PB$, we see that $4OP-76 = OP+19$, so $OP = \frac{95}3$ and the answer is $\boxed{098}$.

 



\textbf{Solution 2}

Reflect triangle $PAM$ across line $AP$, creating an isoceles triangle. Let $x$ be the distance from the top of the circle to point $P$, with $x + 38$ as $AP$. Given the perimeter is 152, subtracting the altitude yields the semiperimeter $s$ of the isoceles triangle, as $114 - x$. The area of the isoceles triangle is:

 $[PAM] = r \cdot s$
$[PAM] = 19 \cdot (114 - x)$

 Now use similarity, draw perpendicular from $O$ to $PM$, name the new point $D$. Triangle $PDO$ is similar to triangle $PAM$, by AA Similarity. Equating the legs, we get:

 $\frac{\sqrt{x}}{19} = \frac{\sqrt{x + 38}}{AM}$

 Solving for $AM$,  it yields $19 \cdot \sqrt{\frac{x + 38}{x}}$.

 $19 \cdot (114 - x) = AM \cdot AP = 19 \cdot (x + 38) \cdot \sqrt{\frac{x + 38}{x}}$

 The $x^3$ cancels, yielding a quadratic. Solving yields $x = \frac{38}{3}$.Add $19$ to find $OP$, yielding $\frac{95}{3}$ or $\boxed{098}$.

 



\textbf{Solution 3}

Let the foot of the perpendicular from $O$ to $PM$ be $D;$ now $OD=19.$ Also let $AM=x$ and $PM=y.$ This means that $OP=\frac{y}{x}\cdot 19$, since $O$ is on the angle bisector of $\angle M.$

 We have that $\tan(\angle AMO)=\frac{19}{x},$ so
 \[\tan(\angle M)=\tan (2\cdot \angle AMO)=\frac{38x}{x^{2}-361}.\]

 However $\tan(\angle M)=\frac{PA}{AM}=\frac{PO+OA}{AM}=\frac{\frac{y}{x}\cdot 19 + 19}{x}$, so
 \[\frac{38x}{x^{2}-361}=19\cdot \frac{\frac{y}{x}+1}{x}\]
 \[\frac{2x^{2}}{x^{2}-361}=\frac{y}{x}+1\]
 \[\frac{x^{2}+361}{x^{2}-361}=\frac{y}{x}.\]
 \[x\cdot \frac{x^{2}+361}{x^{2}-361}=y\]

 We now use the fact that the perimeter of $\triangle PAM$ is $152$:
 \[PO+OA+AM+MP=152\]
 \[\frac{y}{x}\cdot 19+19+x+y=152\]
 \[19\left(\frac{x^{2}+361}{x^{2}-361}\right)+x\cdot \left(\frac{x^{2}+361}{x^{2}-361}\right)+x+19=152\]
 \[(x+19)\left(\frac{x^{2}+361}{x^{2}-361}+\frac{x^{2}-361}{x^{2}-361}\right)=152\]
 \[\frac{2x^{2}}{x-19}=152\]
 \[x^{2}-76x+19\cdot 76=0.\]This quadratic factors as $(x-38)^{2}=0,$ so $x=38$, and 
 \[\frac{y}{x}=\frac{38^{2}+361}{38^{2}-361}=\frac{5}{3}\]
 \[OP=\frac{y}{x}\cdot 19=\frac{95}{3}\to \boxed{98.}\]




\textbf{Solution 4}

Let $K$ be the foot of the altitude from $O$ to $PM,$ and notice how $OK$ is a radius of the circle. Also, we have that $\angle OKM=\angle OAM=90^\circ,$ $OK=OA=90,$ and $OM=OM,$ so $\triangle OKM$ is congruent to $\triangle OAM.$ This means that because $\angle PMO=\angle KMO=\angle AMO,$ $O$ is the foot of the angle bisector from $M.$ 

 Then, let $k=\angle AMO=\angle PMO.$ We have that 
 \[\angle POK=180^\circ-\angle KOM-\angle AOM=2k=\angle PMA,\]so $\triangle PKO$ is similar to $\triangle PAM.$ Let $a=AM=KM,$ and let $p=152$ be the perimeter of $\triangle PAM.$ By similarity, we have that the perimeter of $\triangle PKO$ is $\frac{19}{a}\cdot p,$ so
 \[PK+PO=\dfrac{19}{a}\cdot p-19,\]so
 \[p=PK+PO+OA+AM+MK=\dfrac{19}{a}\cdot p-19+19+2a.\]Solving for $p,$ we have that
 \[p=\dfrac{2a}{1-\frac{19}{a}}\]and using $p=152$ we find that 
 \[a^2-76a+38^2=(a-38)^2=0\]so $a=38.$ 

 Finally, let $b=PO.$ We have by the angle bisector theorem that $PM=2b,$ so using the perimeter information one more time we get $3b+57=152$ and solving gives us
 \[b=\dfrac{95}{3}\implies \boxed{098}.\]




\textbf{Solution 5}

Let $X$ be the point of tangency from the circle to $PM$. Letting $AM=b$, we have $AM=XM=b$. We also know that triangle $POX$ is similar to triangle $PMA$. Taking advantage of the similarity between the sum of the shortest leg and hypotenuse of the two triangles, we have $AP+PM=\frac{b}{19} \cdot (XP+PO)$ $\implies$ $152-b=\frac{b}{19} \cdot (133-2b)$. If we solve this equation, we obtain $b=38$. Now letting  $OP=a$, we know that $PM=2a$, so $3a+57=152$, or $OP=a= \frac{95}{3}$, so our desired answer is $\boxed{098}$.~Irfans123

 


	\item (\textit{AIME II 2002}) Circles $\mathcal{C}_{1}$ and $\mathcal{C}_{2}$ intersect at two points, one of which is $(9,6)$, and the product of the radii is $68$. The x-axis and the line $y = mx$, where $m > 0$, are tangent to both circles. It is given that $m$ can be written in the form $a\sqrt {b}/c$, where $a$, $b$, and $c$ are positive integers, $b$ is not divisible by the square of any prime, and $a$ and $c$ are relatively prime. Find $a + b + c$.


 
 \textbf{Answer: }282



\textbf{Solution 1}

Let the smaller angle between the $x$-axis and the line $y=mx$ be $\theta$. Note that the centers of the two circles lie on the angle bisector of the angle between the $x$-axis and the line $y=mx$. Also note that if $(x,y)$ is on said angle bisector, we have that $\frac{y}{x}=\tan{\frac{\theta}{2}}$.  Let $\tan{\frac{\theta}{2}}=m_1$, for convenience. Therefore if $(x,y)$ is on the angle bisector, then $x=\frac{y}{m_1}$. Now let the centers of the two relevant circles be 

 


\end{enumerate}